\subsection{下一阶段四周安排}

\small
\begin{longtable}{p{1.3cm}p{1.3cm}p{3.9cm}p{4.7cm}p{4.15cm}}
\caption{四周训练任务与核对办法}\label{tab:four-week-plan}\\
\toprule
时间 & 负责人 & 任务 & 产物 & 核对办法\\
\midrule
\endfirsthead
\toprule
时间 & 负责人 & 任务 & 产物 & 核对办法\\
\midrule
\endhead
\bottomrule
\endfoot
第1周 & 队长 & 用一道A题完成ODE/RK4基线，并完成一个二维扩散FDM算例。
& 方程说明、两套脚本、三档步长或网格结果。 & 量纲正确；误差随加密下降；另一名成员可按入口复算一张表。\\
第1周 & 组员1 & 用小规模排程题对照枚举与MILP。 & 变量表、约束表、枚举结果和求解日志。 & 小样例目标值一致；每条约束能指出题意来源。\\
第1周 & 组员2 & 完成一份附件数据的字段、缺失、异常和泄漏检查。 & 数据字典、清洗记录、训练/验证划分脚本。 & 每次删除或填补都有对象和理由；预处理不读取验证集统计量。\\
第2周 & 组员1 & 在同一连续优化题上比较遍历、PSO和NSGA--II的对应任务。 & 参数配置、20个种子结果、收敛与可行率表。 & 计算预算一致；报告中位数和离散程度，不只报最好一次。\\
第2周 & 组员2 & 完成PCA--聚类和熵权--TOPSIS两个小例。 & 解释率、载荷、簇中心、权重扰动与排序表。 & 能说明主成分、簇和排名分别回答什么问题。\\
第3周 & 全队 & 选择一题做36小时压缩训练。 & 可编译论文、代码入口、结果目录和复盘单。 & 正文中每个关键数值可追到结果文件；无未定义引用。\\
第4周 & 全队 & 安排72小时完整训练并互换审稿。 & 论文、源代码、附件、个人问题清单。 & 12小时内定题，24小时给基线，48小时冻结主结果，交稿前完成复算。\\
\end{longtable}
\normalsize

